\chapter{光的吸收、散射和色散}

\problem{6.1}{一固体有两个吸收带，宽度都是 $30\,\mathrm{nm}$。一带处在蓝光区（$450\,\mathrm{nm}$ 附近），另一带处在黄光区（$580\,\mathrm{nm}$ 附近）。设第一带吸收系数为 $50\,\mathrm{cm^{-1}}$，第二带的为 $250\,\mathrm{cm^{-1}}$。试绘出白光分别透过 $0.1\,\mathrm{mm}$ 及 $5\,\mathrm{mm}$ 的该物质后在吸收带附近光强分布的概况。}
\answerarea

\problem{6.2}{某种介质的吸收系数 $\alpha_{\mathrm{a}}$ 为 $0.32\,\mathrm{cm^{-1}}$，求透射光强为入射光强的 $0.1$、$0.2$、$0.5$ 及 $0.8$ 倍时，该介质的厚度各为多少？}
\answerarea
\clearpage

\problem{6.3}{如果同时考虑到吸收和散射都将使透射光强度减弱，则透射光表达式中的 $\alpha$ 可看做是由两部分合成，一部分 $\alpha_{\mathrm{a}}$ 是由于真正的吸收（变为物质分子的热运动），另一部分 $\alpha_{\mathrm{s}}$（称为散射系数）是由于散射，于是该式可写作 $I=I_0\mathrm{e}^{-(\alpha_{\mathrm{a}}+\alpha_{\mathrm{s}})d}$。如果光通过一定厚度的某种物质后，只有 $20\%$ 的光强通过，已知该物质的散射系数等于吸收系数的 $1/2$。假定不考虑散射，则透射光强可增加多少？}
\answerarea

\problem{6.4}{计算波长为 $253.6\,\mathrm{nm}$ 和 $546.1\,\mathrm{nm}$ 的两条谱线瑞利散射的强度之比。}
\answerarea
\clearpage

\problem{6.5}{太阳光束由小孔射入暗室，室内的人沿着与光束垂直及与之成 $45^\circ$ 的方向观察这光束时，见到瑞利散射的散射光强度之比等于多少？}
\answerarea

\problem{6.6}{一束光通过液体，用尼科耳棱镜正对这束光进行观察，当尼科耳棱镜主截面竖直时，光强达最大值；当尼科耳棱镜主截面水平时，光强为零。再从侧面观察其散射光，在尼科耳棱镜主截面为竖直和水平两个位置时，光强之比为 $20:1$，计算散射光的退偏振度。}
\answerarea
\clearpage

\problem{6.7}{一种光学玻璃对于波长 $435.8\,\mathrm{nm}$ 和 $546.1\,\mathrm{nm}$ 的折射率分别为 $1.613\,0$ 和 $1.602\,6$，试应用柯西公式来计算这种玻璃对波长 $600\,\mathrm{nm}$ 的光的色散 $\dfrac{\mathrm{d}n}{\mathrm{d}\lambda}$。}
\answerarea

\problem{6.8}{一种光学玻璃对汞蓝光 $435.8\,\mathrm{nm}$ 和汞绿光 $546.1\,\mathrm{nm}$ 的折射率分别为 $1.652\,50$ 和 $1.624\,50$。（1）试用柯西公式计算公式中的常量 $a$ 和 $b$；（2）并求它对 $589\,\mathrm{nm}$ 钠黄光的折射率 $n$ 和色散 $\dfrac{\mathrm{d}n}{\mathrm{d}\lambda}$。}
\answerarea
\clearpage

\problem{6.9}{一个顶角为 $60^\circ$ 的棱镜由某种玻璃制成，它的色散特性可用柯西公式中的常量 $a=1.416$，$b=1.72\times 10^{-10}\,\mathrm{cm^2}$ 来表示。将棱镜的位置放置得使它对 $600\,\mathrm{nm}$ 的波长产生最小偏向角，试计算这个棱镜的角色散率为多少。}
\answerarea

\problem{6.10}{波长为 $0.67\,\mathrm{nm}$ 的 X 射线，由真空入射到某种玻璃时，在掠射角不超过 $0.1^\circ$ 的条件下发生全反射，试计算玻璃对这波长的折射率，并解释所得的结果。}
\answerarea
\clearpage

\sproblem{6.11}{波长为 $0.07\,\mathrm{nm}$ 的 X 射线，其折射率比 $1$ 小 $1.60\times 10^{-6}$，试求全反射时，X 射线的掠射角为多大？}
\answerarea

\sproblem{6.12}{对于波长 $\lambda=400\,\mathrm{nm}$ 的光，某种玻璃的折射率 $n=1.63$，对 $\lambda'=500\,\mathrm{nm}$ 的光，其折射率 $n'=1.58$。假如柯西公式的近似形式 $n=a+\dfrac{b}{\lambda^2}$ 适用，试求此种玻璃对波长为 $600\,\mathrm{nm}$ 的光的色散 $\dfrac{\mathrm{d}n}{\mathrm{d}\lambda}$ 值。}
\answerarea
\clearpage

\sproblem{6.13}{研究性课题：$\mathrm{PM2.5}$ 与大气环境污染检测。}
\answerarea
\clearpage
